\subsection{K0 holds in every partition monoid}\label{subsec:K0}

\emph{(16 Oct 2026, evening; files \texttt{pa8/k0/}.)} The second order-$8$ condition is shorter to prove than K4, because its $z$-conditions make $z$ a common matched partner of $a$ and $b$, which pins every block of $z$ to a ``cell'', and the $y$-conditions then force the connection of two $\T(a)$-blocks to run through the diagonal cells.

\begin{theorem}\label{thm:K0}
Let $a,b,c,y,z$ be projections of $\P_X$ with $a\bowtie c$, $c\bowtie b$, $yc=cy$, $\theta_y(b)=\theta_b(a)$, $\theta_b(z)=b$, $\theta_z(a)=\theta_z(b)$ and $\theta_y(z)=\theta_a(b)$. Then $a\bowtie b$. Consequently \textup{(E$_1$)} holds in every $\P_X$ (Theorem~\ref{thm:K0E1}, with (W$_2$) from Theorem~\ref{thm:Wn}), $\mathfrak P_8^1$ embeds in no $\P(\P_X)$, and $\mathfrak P_8^1$ is a minimal non-embeddable projection algebra (as in Theorem~\ref{thm:K4}: every proper subalgebra embeds, Section~\ref{sec:comp}). With Theorem~\ref{thm:K4}, all seven candidates of order $8$ are proved minimal non-embeddable.
\end{theorem}

\begin{proof}
We use the notation of Theorem~\ref{thm:R}: $A_i,C_i,B_i$ ($1\le i\le m$) are the $\T$-blocks of $a,c,b$, indexed so that $\Gamma_i\ni A_i,C_i$ are the classes of $\varepsilon_a\vee\varepsilon_c$ meeting $\supp a\cup\supp c$ and $\Delta_i\ni C_i,B_i$ those of $\varepsilon_b\vee\varepsilon_c$ (Lemma~\ref{lem:PX}(c)); in particular every block of $c$ meeting $A_i$ lies in $\Gamma_i$, every block of $b$ meeting $C_i$ lies in $\Delta_i$, and no $\T$-block of one of $a,b,c$ is an $\N$-block of a matched partner. Put $u=\theta_a(b)$, $v=\theta_b(a)$, $n=\rk u$. The blocks of $u$ are the $\N(a)$-blocks and the sets $M\cap\supp a$ for the classes $M$ of $\varepsilon_a\vee\varepsilon_b$ meeting $\supp a$, transversal iff $M$ meets $\supp b$ (Lemma~\ref{lem:PX}(a)); so $a\bowtie b$ iff $n=m$ and every class of $\varepsilon_a\vee\varepsilon_b$ meeting $\supp a$ meets $\supp b$ (Lemma~\ref{lem:PX}(c) and its proof).

\smallskip\emph{Step 1: algebra, and the cells.} $u\le a$, $v\le b$, and $\theta_y(u)=\theta_y(\theta_y(z))=\theta_y(z)=u$, $\theta_y(v)=\theta_y(\theta_y(b))=v$ give $u\le y$, $v\le y$; by \eqref{eq:Q1}, $ay=ya$ and $by=yb$, and then $\theta_y(a)=ya=u$, $\theta_y(b)=yb=v$. Put $p=\theta_z(a)=\theta_z(b)$. Then $b=\theta_b(z)\bowtie\theta_z(b)=p$ (Lemma~\ref{lem:friends}), so $\rk p=m$. The $\T$-blocks of $p=\theta_z(a)$ are the sets $L\cap\supp z$ for the classes $L$ of $\varepsilon_a\vee\varepsilon_z$ meeting both $\supp a$ and $\supp z$; there are $m$ of them, each contains a $\T(a)$-block, and there are $m$ $\T(a)$-blocks; so every class of $\varepsilon_a\vee\varepsilon_z$ meeting $\supp a$ meets $\supp z$ and contains exactly one $\T(a)$-block, which by Lemma~\ref{lem:PX}(a) says $\theta_a(z)=a$ (hence also $a\bowtie p$, by Lemma~\ref{lem:friends}, though we shall only use $\theta_a(z)=a$ and $\theta_b(z)=b$). Let $L_i$ be the class of $\varepsilon_a\vee\varepsilon_z$ containing $A_i$ and $L'_i$ the class of $\varepsilon_b\vee\varepsilon_z$ containing $B_i$; the $L_i$ are distinct, and so are the $L'_i$ (this is $\theta_b(z)=b$). A \emph{cell} is a nonempty set $L\cap L'$ with $L$ a class of $\varepsilon_a\vee\varepsilon_z$ and $L'$ a class of $\varepsilon_b\vee\varepsilon_z$; the cells partition $X$, every block of $z$ lies in a cell, every block of $a$ lies in one class of $\varepsilon_a\vee\varepsilon_z$, and every block of $b$ in one class of $\varepsilon_b\vee\varepsilon_z$. The \emph{diagonal} cells are $D_i=L_i\cap L'_i$ (possibly empty).

\smallskip\emph{Step 2: $y$ against $a,b,c$.} Lemma~\ref{lem:PX}(b) applies to the commuting pairs $(y,a)$, $(y,b)$, $(y,c)$. We claim that for $(y,a)$ and $(y,b)$ there is no class of type (iii): \emph{every $\T(a)$-block and every $\T(b)$-block meets $\supp y$.} Let $N$ be an $\N(y)$-block which is a union of $\T(a)$-blocks, $N\supseteq A_i$. As an $\N(y)$-block, $N$ is a nontransversal block of $\theta_y(b)=v=\theta_b(a)$, i.e.\ an $\N(b)$-block or a set $M\cap\supp b$ for a class $M$ of $\varepsilon_a\vee\varepsilon_b$ missing $\supp a$; the latter is impossible as $N\supseteq A_i$, so $N$ is an $\N(b)$-block. For $(y,c)$, $N$ is of type (iii) or (iv). If $N$ is an $\N(c)$-block, it meets $A_i$, so $N\subseteq\Gamma_i$, and $N\subseteq\supp a$ gives $N=A_i$, a $\T(a)$-block that is an $\N(c)$-block, impossible. So $N$ is a union of $\T(c)$-blocks, and the $\N(b)$-block $N$ contains some $C_k$, impossible. Next let $N$ be an $\N(y)$-block which is a union of $\T(b)$-blocks, $N\supseteq B_k$. It is a nontransversal block of $v=\theta_b(a)$ which is not an $\N(b)$-block, so $N=M\cap\supp b$ for a class $M$ of $\varepsilon_a\vee\varepsilon_b$ missing $\supp a$; thus $N$ contains no $\T(a)$-block, and by what was just proved (no type (iii) for $(y,a)$) $N$ is an $\N(a)$-block. For $(y,c)$, if $N$ were an $\N(c)$-block then $N\subseteq\Delta_k$ and $N=B_k$, impossible; so $N$ is a union of $\T(c)$-blocks. Every $\T(b)$- and every $\T(c)$-block lies in some $\Delta_l$, so $N=\bigsqcup_l(N\cap\Delta_l)$, where $N\cap\Delta_l$ is both a union of $\T(b)$-blocks (hence $B_l$ or empty) and a union of $\T(c)$-blocks (hence $C_l$ or empty): $N$ is a disjoint union of blocks $B_l=C_l$. Being a block of $a$, $N$ lies in one $\Gamma_l$, so $N=B_k=C_k$ is an $\N(a)$-block which is a $\T(c)$-block, impossible.

Consequences. Every $\N(y)$-block is an $\N(a)$-block and an $\N(b)$-block (type (iv) for $(y,a)$ and $(y,b)$); every $\T(y)$-block $Y$ lies in a class of type (i) or (ii), so $Y$ lies in a union of $\T(a)$-blocks or in an $\N(a)$-block, in a union of $\T(b)$-blocks or in an $\N(b)$-block, and in a union of $\T(c)$-blocks or in an $\N(c)$-block; and every $\N$-block of $a$, $b$ or $c$ meeting a $\T(y)$-block contains it (types (ii) and (iv)). The $\T$-blocks of $u=ya$ are the classes $U_1,\dots,U_n$ of type (i) of $\varepsilon_y\vee\varepsilon_a$, each a union of $\T(a)$-blocks and of $\T(y)$-blocks, and its $\N$-blocks are the $\N(a)$-blocks; comparing with the description of $u=\theta_a(b)$ above, the $U_j$ are the sets $M_j\cap\supp a$ for the classes $M_1,\dots,M_n$ of $\varepsilon_a\vee\varepsilon_b$ meeting $\supp a$, and every $M_j$ meets $\supp b$. Put $I_j=\{i:A_i\subseteq U_j\}$; the $I_j$ partition $[m]$, and $a\bowtie b$ iff $n=m$ iff every $|I_j|=1$.

\smallskip\emph{Step 3: the classes of $\varepsilon_y\vee\varepsilon_z$.} By Lemma~\ref{lem:PX}(a) the blocks of $\theta_y(z)$ inside $\supp y$ are the sets $K\cap\supp y$ for the classes $K$ of $\varepsilon_y\vee\varepsilon_z$ meeting $\supp y$, transversal iff $K$ meets $\supp z$; the blocks of $u=ya$ inside $\supp y$ are the $U_j$ (transversal) and the $\N(a)$-blocks contained in $\supp y$ (nontransversal; type (ii)). As $\theta_y(z)=u$, the classes of $\varepsilon_y\vee\varepsilon_z$ meeting $\supp y$ are $K_1,\dots,K_n$ with $K_j\cap\supp y=U_j$ and $K_j\cap\supp z\ne\emptyset$, and, for each $\N(a)$-block $F\subseteq\supp y$, a class $K_F$ with $K_F\cap\supp y=F$. In particular each $A_i$ lies in $K_j$ for $i\in I_j$, and each $\N(a)$-block is an $\N(y)$-block or lies in its $K_F$.

\smallskip\emph{Step 4: blocks of $y$ and cells.} Call a $\T(y)$-block \emph{crossing} if it meets two $\T(a)$-blocks.

(a) \emph{A crossing block $Y$ satisfies $Y\subseteq\bigcup_i(A_i\cap B_i\cap C_i)$, the union over the $i$ with $Y\cap A_i\ne\emptyset$; in particular $Y$ lies in the union of the diagonal cells.} $Y$ meets $\supp a$, so it lies in a type (i) class $U_j$ and $Y\subseteq\supp a$. $Y$ is not inside an $\N(c)$-block $F'$: such an $F'$ would meet $A_i$ and $A_{i'}$ for two $i\ne i'$, so $F'\subseteq\Gamma_i\cap\Gamma_{i'}$. So $Y\subseteq\bigcup_kC_k$; a point of $Y\cap A_i$ lies in some $C_k\subseteq\Gamma_k$ and in $A_i\subseteq\Gamma_i$, so $k=i$: $Y\cap A_i\subseteq C_i$. A point $x\in Y\cap C_i$ lies in a block of $b$ inside $\Delta_i$, namely $B_i$ or an $\N(b)$-block $G\subseteq\Delta_i$; in the second case $G\supseteq Y$, but $Y\cap A_{i'}\subseteq C_{i'}\subseteq\Delta_{i'}$ for a second index $i'$. So $Y\cap A_i\subseteq A_i\cap B_i\cap C_i$, and a point of $A_i\cap B_i$ lies in $L_i\cap L'_i=D_i$.

(b) \emph{A non-crossing $\T(y)$-block $Y$ lies in one cell.} If $Y$ meets no $\T(a)$-block then $Y$ lies in an $\N(a)$-block, and if it meets exactly one, $A_i$, then it lies in $U_j\subseteq\supp a$ and $Y\subseteq A_i$; either way $Y$ lies in one block of $a$, hence in one class of $\varepsilon_a\vee\varepsilon_z$ and in one $\Gamma$-class or $\T$-less class of $\varepsilon_a\vee\varepsilon_c$. For $\varepsilon_b\vee\varepsilon_z$: if $Y$ lies in an $\N(b)$-block we are done; otherwise $Y\subseteq\bigcup_kB_k$, and we claim $Y$ meets only one $B_k$. Suppose $Y$ meets $B_k$ and $B_l$, $k\ne l$. $Y$ is not inside an $\N(c)$-block (it would lie in $\Delta_k\cap\Delta_l$), so $Y\subseteq\bigcup C$'s, and as in (a) $Y\cap B_k\subseteq C_k\subseteq\Gamma_k$, $Y\cap B_l\subseteq C_l\subseteq\Gamma_l$, contradicting the fact that $Y$ lies in one class of $\varepsilon_a\vee\varepsilon_c$. So $Y\subseteq B_k\subseteq L'_k$.

(c) \emph{Every $\N(y)$-block and every block of $z$ lies in one cell}: an $\N(y)$-block is a block of $a$ and of $b$ (Step~2).

(d) \emph{Every class $K$ of $\varepsilon_y\vee\varepsilon_z$ containing no crossing block lies in one cell, and every class containing a crossing block lies in $\bigcup_iD_i$.} $K$ is connected by chains of blocks of $y$ and $z$ in which consecutive blocks meet. By (b), (c) a non-crossing block lies in one cell, so it does not leave the cell of the block before it; by (a) a crossing block lies in $\bigcup D_i$, and a non-crossing block meeting it lies in some $D_i$. Starting the chains at a crossing block, resp.\ noting that without crossing blocks no chain changes cell, gives the claim.

\smallskip\emph{Step 5: $\varepsilon_a\le\varepsilon_b\vee\varepsilon_z$.} Let $A_i\subseteq U_j\subseteq K_j$. If $K_j$ contains no crossing block, $A_i$ lies in one cell. Otherwise $K_j\subseteq\bigcup_{i'}D_{i'}$, and $A_i\subseteq L_i$ with $D_{i'}\subseteq L_{i'}$ and the $L$'s distinct, so $A_i\subseteq D_i\subseteq L'_i$. An $\N(a)$-block $F\subseteq\supp y$ lies in $K_F$, which meets $\supp a$ in nothing (its intersection with $\supp y$ is $F$) and so contains no crossing block (crossing blocks lie in $\supp a$); so $F$ lies in one cell. An $\N(a)$-block not inside $\supp y$ is an $\N(y)$-block, in one cell by Step~4(c). So every block of $a$ lies in one class of $\varepsilon_b\vee\varepsilon_z$, and therefore $\varepsilon_a\vee\varepsilon_b\le\varepsilon_b\vee\varepsilon_z$.

\smallskip\emph{Step 6: the contradiction.} Suppose $n<m$, so that some $I_j$ contains two indices, say $1,2\in I_j$. The class $U_j$ of type (i) is connected by chains alternating $\T(a)$- and $\T(y)$-blocks, so some $\T(y)$-block of $U_j$ meets $A_1$ and $A_2$: $K_j$ contains a crossing block, and by Step~5, $A_1\subseteq L'_1$ and $A_2\subseteq L'_2$, two distinct classes of $\varepsilon_b\vee\varepsilon_z$. But $A_1,A_2\subseteq M_j$, a single class of $\varepsilon_a\vee\varepsilon_b$, which by Step~5 lies in one class of $\varepsilon_b\vee\varepsilon_z$. Hence $n=m$, every class of $\varepsilon_a\vee\varepsilon_b$ meeting $\supp a$ meets $\supp b$ (Step~2), and $a\bowtie b$.
\end{proof}

Every step was tested by the dropped-hypothesis map before being written down (\texttt{pa8/k0/k0claims.py}: all instances of $\P(\P_4)$, and a sample of $\P(\P_5)$; the statements are: $\theta_a(z)=a$; $\N(y)$-blocks are blocks of $a,b,c$; crossing blocks lie in $\bigcup(A_i\cap B_i\cap C_i)$; non-crossing $\T(y)$-blocks and $\N(y)$-blocks lie in one cell; classes of $\varepsilon_y\vee\varepsilon_z$ lie in one cell or in the diagonal; $\varepsilon_a\le\varepsilon_b\vee\varepsilon_z$). The map shows where each hypothesis enters: $\theta_b(z)=b$ and $\theta_z(a)=\theta_z(b)$ are used only through Step~1, i.e.\ through $\theta_a(z)=a$ and $\theta_b(z)=b$ (the $L_i$ distinct, the $L'_i$ distinct) --- so the proof gives also the condition with $\theta_z(a)=\theta_z(b)$ replaced by $\theta_a(z)=a$; $\theta_y(z)=\theta_a(b)$ gives $y\sim a$, $y\sim b$ (Step~1) and the classes $K_j$, $K_F$ (Step~3), without which Step~4(a) and Step~5 fail; $yc=cy$ and $\theta_y(b)=\theta_b(a)$ carry Step~2 and the confinement of the blocks of $y$ in Step~4 (without either, non-crossing blocks leave their cells). The proof is shorter than that of K4 because here $z$ is matched with both $a$ and $b$ (Step~1), so that its blocks are confined to cells from the start, whereas in K4 the roles of $y$ and $z$ are mixed ($z$ sees $y$ as $v$ and $a$ as $u$) and the confinement has to be established along chains (Step~7 there). Remarkably, the whole argument never uses the $\T$-blocks of $z$: only the classes of $\varepsilon_a\vee\varepsilon_z$ and $\varepsilon_b\vee\varepsilon_z$ enter, through the cells.

